\ncdmodultestmfour{1}{Arbejdsliv, Tid og Fleksibilitet}{全真模考卷一：四天工作制、职场转行与工作生活平衡}{四天工作制完形 / 银行顾问转型木匠补句 / 兼职护士开放完形 / 年假·A-kasse·产假·图书馆等 6 则 / 回信 Mikkel 与居家办公辩论 / 工作生活与育儿平权 Emneopgave 及 Arbejde og fritid 口语讨论}{\textit{-s} 与 \textit{blive} 被动 · 假拟条件句 · 固定介词搭配（\textit{mangel på, nødt til, ret til}）· 情态小品词 · \textit{Debatindlæg} 让步反驳（\textit{Ganske vist..., men...}）}

% =====================================================================
% DEL 1: LÆSNING (50 minutter · 29 point)
% =====================================================================
\ncdsubsec{Del 1: Læsning（阅读理解｜50 minutter · 29 point）}

\begin{tipbox}{考场规则与配速建议（Læsning｜50 minutter · 禁用词典）}
\small
Læsning Opgave 1, 2, 3 og 4 udleveres samtidig. Samlet tid: \textbf{50 minutter}. Kursisten må \textbf{ikke} bruge hjælpemidler（禁用任何词典、笔记或手机）. Der gives 1 point for hvert korrekt svar (maksimum \textbf{29 point}: 8 + 5 + 10 + 6). \textit{Bestået / Ikke bestået} 及格线由各语言中心依 SIRI 指引划定（通常约 60--65\%，建议自测稳定达到 \textbf{20+ point}）。建议配速：Opgave 1（10 min）$\rightarrow$ Opgave 2（12 min）$\rightarrow$ Opgave 3（15 min）$\rightarrow$ Opgave 4（10 min），最后留 3 分钟复查。
\end{tipbox}

% ---------------------------------------------------------------------
\subsubsection*{Opgave 1 (8 point): \textit{Fire dages arbejdsuge i praksis}}

\begin{emailbox}{Kursistark --- Instruktion（Opgave 1）}
\begin{itemize}
  \item Læs teksten \textit{Fire dages arbejdsuge i praksis}.
  \item I teksten mangler der otte ord/udtryk (1--8).
  \item I skemaet er der fire ord/udtryk: A, B, C og D.
  \item Sæt kryds i skemaet ved det ord/udtryk, der passer til hvert af hullerne i teksten. Se eksemplet (0).
\end{itemize}

\vspace{0.35em}
\hrule
\vspace{0.45em}

\textbf{\color{NCDNavy}Fire dages arbejdsuge i praksis}

\vspace{0.3em}
De fleste danskere arbejder stadig omkring 37 timer fordelt på fem hverdage, men i de senere år har flere virksomheder og enkelte kommuner afprøvet en arbejdsuge på kun fire dage. Ideen lyder måske \textbf{(0)}\ \ldots, fordi medarbejderne får en ekstra fridag uden at gå ned i løn.

\vspace{0.25em}
Erfaringerne fra forsøgene viser faktisk en række tydelige \textbf{(1)}\ \ldots\ ved den kortere arbejdsuge. Når personalet ved, at fredagen er fri, bliver mange bedre til at prioritere de vigtigste opgaver og holde kortere møder. Flere ledere fortæller, at medarbejderne møder ind med mere energi og er mindre \textbf{(2)}\ \ldots\ efter weekenden. Samtidig falder sygefraværet ofte, fordi der bliver bedre tid til både familie og hvile.

\vspace{0.25em}
Alligevel er en fire dages arbejdsuge ikke lige nem at \textbf{(3)}\ \ldots\ på alle arbejdspladser. På et kontor eller i en IT-virksomhed kan man ofte arbejde mere effektivt og nå det samme på fire dage. \textbf{(4)}\ \ldots\ er det langt sværere på et hospital, i en børnehave eller i hjemmeplejen, hvor der hele tiden skal være mennesker til stede for at passe på borgerne.

\vspace{0.25em}
En anden \textbf{(5)}\ \ldots\ er arbejdspresset på de fire arbejdsdage. Hvis virksomheden ikke reducerer mængden af opgaver, kan de fire dage hurtigt blive alt for travle og stressede. Nogle ansatte oplever, at de er \textbf{(6)}\ \ldots\ at arbejde længere om aftenen eller svare på mails i deres fritid for at nå det hele.

\vspace{0.25em}
\textbf{(7)}\ \ldots\ peger flere arbejdsmarkedsforskere på, at fleksibilitet ikke ser ens ud for alle brancher. For nogle fungerer en tredages weekend bedst, mens andre hellere vil arbejde kortere dage eller have mulighed for hjemmearbejde. \textbf{(8)}\ \ldots\ ser fire dages arbejdsuge ud til at blive et fast tilbud i stadig flere moderne virksomheder, så længe opgaverne tilpasses realistisk.
\end{emailbox}

\small
\renewcommand{\arraystretch}{1.28}
\begin{longtable}{@{\extracolsep{\fill}} >{\bfseries\color{NCDNavy}\centering\arraybackslash}p{0.8cm} >{\raggedright\arraybackslash}p{3.45cm} >{\raggedright\arraybackslash}p{3.45cm} >{\raggedright\arraybackslash}p{3.45cm} >{\raggedright\arraybackslash}p{3.45cm} @{}}
\toprule
\textbf{Nr.} & \textbf{A} & \textbf{B} & \textbf{C} & \textbf{D} \\
\midrule
\endfirsthead
\toprule
\textbf{Nr.} & \textbf{A} & \textbf{B} & \textbf{C} & \textbf{D} \\
\midrule
\endhead
\bottomrule
\endfoot
\textcolor{NCDGray}{0} & \textcolor{NCDGray}{kedelig} \hfill \fbox{\phantom{X}} & \textcolor{NCDGray}{besværlig} \hfill \fbox{\phantom{X}} & \textcolor{NCDGray}{fristende} \hfill \fbox{\textbf{X}} & \textcolor{NCDGray}{tilfældig} \hfill \fbox{\phantom{X}} \\
\midrule
1 & ulemper \hfill \fbox{\phantom{X}} & gevinster \hfill \fbox{\phantom{X}} & omkostninger \hfill \fbox{\phantom{X}} & begrænsninger \hfill \fbox{\phantom{X}} \\
\midrule
2 & udkørte \hfill \fbox{\phantom{X}} & stolte \hfill \fbox{\phantom{X}} & engagerede \hfill \fbox{\phantom{X}} & forberedte \hfill \fbox{\phantom{X}} \\
\midrule
3 & aflyse \hfill \fbox{\phantom{X}} & opgive \hfill \fbox{\phantom{X}} & forbyde \hfill \fbox{\phantom{X}} & indføre \hfill \fbox{\phantom{X}} \\
\midrule
4 & Derfor \hfill \fbox{\phantom{X}} & Desuden \hfill \fbox{\phantom{X}} & Derimod \hfill \fbox{\phantom{X}} & Nemlig \hfill \fbox{\phantom{X}} \\
\midrule
5 & udfordring \hfill \fbox{\phantom{X}} & mulighed \hfill \fbox{\phantom{X}} & løsning \hfill \fbox{\phantom{X}} & erfaring \hfill \fbox{\phantom{X}} \\
\midrule
6 & bange for \hfill \fbox{\phantom{X}} & glade for \hfill \fbox{\phantom{X}} & stolte af \hfill \fbox{\phantom{X}} & nødt til \hfill \fbox{\phantom{X}} \\
\midrule
7 & Ellers \hfill \fbox{\phantom{X}} & Derudover \hfill \fbox{\phantom{X}} & Desværre \hfill \fbox{\phantom{X}} & Pludselig \hfill \fbox{\phantom{X}} \\
\midrule
8 & For eksempel \hfill \fbox{\phantom{X}} & Tværtimod \hfill \fbox{\phantom{X}} & Alt i alt \hfill \fbox{\phantom{X}} & Desværre \hfill \fbox{\phantom{X}} \\
\end{longtable}
\normalsize

% ---------------------------------------------------------------------
\subsubsection*{Opgave 2 (5 point): \textit{Fra bankrådgiver til snedker}}

\begin{emailbox}{Kursistark --- Instruktion（Opgave 2）}
\begin{itemize}
  \item Læs teksten \textit{Fra bankrådgiver til snedker}.
  \item I hvert afsnit er der et hul, hvor der mangler en sætning.
  \item Sætningerne, der mangler, står på listen (A--H).
  \item Find den sætning, der passer bedst i hvert afsnit (1--5).
  \item Skriv dit svar (= bogstav) på linjen i hullet.
  \item Der er \textbf{to sætninger, du ikke skal bruge}. Se eksemplet i første afsnit.
\end{itemize}

\vspace{0.35em}
\hrule
\vspace{0.45em}

\textbf{\color{NCDNavy}Fra bankrådgiver til snedker}

\vspace{0.3em}
Da Anders Holm var i slutningen af tyverne, så hans fremtid helt fastlagt ud. Han havde taget en solid uddannelse inden for økonomi og arbejdet fem år som bankrådgiver i Aarhus. Lønnen var god, og kollegaerne var flinke, men alligevel voksede en stille uro indeni ham, hver gang han så ud ad vinduet mod gågaden. \textbf{(0)}\ \underline{\hspace{0.45cm}\textbf{A}\hspace{0.45cm}}\textbf{.} Hver aften lukkede han sin bærbare computer med en følelse af, at han slet ikke havde skabt noget konkret i løbet af dagen.

\vspace{0.25em}
Vendepunktet kom en regnfuld weekend i november, hvor Anders meldte sig til et kort aftenkursus i møbelsnedkeri på den lokale aftenskole. Egentlig ville han bare lære at reparere et gammelt spisebord til sin lejlighed. \textbf{(1)}\ \rule{1.8cm}{0.4pt}. Den særlige duft af friskhøvlet eg og glæden ved at forme træ med sine egne hænder ramte ham øjeblikkeligt. Allerede efter den tredje aften begyndte han at låne fagbøger om træsorter, samlinger og værktøj på biblioteket.

\vspace{0.25em}
Efter et halvt år tog Anders en modig beslutning, som overraskede både hans familie og venner. Han sagde sin trygge stilling i banken op og begyndte på en erhvervsuddannelse som snedker. \textbf{(2)}\ \rule{1.8cm}{0.4pt}. Som voksenlærling tjente han pludselig langt mindre end før, og derfor måtte han og kæresten lægge deres budget helt om og droppe dyre udlandsrejser i et par år.

\vspace{0.25em}
Heldigvis fandt Anders hurtigt en læreplads på et mindre, lokalt værksted uden for Odder. Her bestod hverdagen af både fint præcisionsarbejde og tunge løft, når der skulle bygges specialdesignede køkkener og reoler til private kunder. \textbf{(3)}\ \rule{1.8cm}{0.4pt}. Samtidig var værkstedets mester glad for en lærling, der havde styr på økonomi, tilbudsgivning og kundekontakt fra sit tidligere arbejdsliv.

\vspace{0.25em}
I dag, fire år efter sit store karriereskift, er Anders udlært møbelsnedker og har åbnet sit eget værksted sammen med en tidligere studiekammerat. De fokuserer på holdbare møbler lavet af dansk træ fra bæredygtige skove. \textbf{(4)}\ \rule{1.8cm}{0.4pt}. Kunderne vil nemlig gerne betale lidt ekstra, når de kan tale direkte med håndværkeren og se præcis, hvor materialerne kommer fra.

\vspace{0.25em}
Når Anders ser tilbage på sine år i banken, fortryder han hverken uddannelsen eller forandringen. Skiftet har lært ham, at succes ikke kun handler om lønseddel og status, men om at glæde sig til at møde på arbejde mandag morgen. \textbf{(5)}\ \rule{1.8cm}{0.4pt}. ”Det vigtigste er, at man tør lytte til sig selv, før hverdagen føles som en vane uden retning,” siger han med et smil.
\end{emailbox}

\begin{grammatikbox}{Liste (A--H)（A 为例句 · 5 项作答 + 2 项多余干扰项）}
\small
\begin{itemize}[leftmargin=2.2em, itemsep=0.35em]
  \item[\textbf{\color{NCDGray}A.}] \textcolor{NCDGray}{Han tilbragte næsten hele sin arbejdsdag foran en skærm med regneark, møder og tal.} {\footnotesize\color{NCDRed}[Eksempel (0)]}
  \item[\textbf{B.}] Alligevel savner han stadig de store bonusser og det høje tempo i finansverdenen.
  \item[\textbf{C.}] I begyndelsen var hans ryg og hænder ømme hver eftermiddag, fordi kroppen skulle vænne sig til det fysiske arbejde.
  \item[\textbf{D.}] Derfor opfordrer han ofte andre til at overveje et sporskifte, hvis de føler sig fastlåst i deres job.
  \item[\textbf{E.}] Men oplevelsen i værkstedet tændte en uventet begejstring hos ham.
  \item[\textbf{F.}] Desuden tilbyder flere banker nu billige lån til unge, der vil købe deres første hus på landet.
  \item[\textbf{G.}] Overgangen fra fast bankløn til elevløn krævede dog en del tilpasning i privatøkonomien.
  \item[\textbf{H.}] Selvom konkurrencen fra store møbelkæder er hård, har det lille firma allerede travlt med bestillinger.
\end{itemize}
\end{grammatikbox}

% ---------------------------------------------------------------------
\subsubsection*{Opgave 3 (10 point): \textit{Sofie arbejder som sygeplejerske på deltid}（无词库开放完形）}

\begin{emailbox}{Kursistark --- Instruktion（Opgave 3）}
\begin{itemize}
  \item Læs teksten \textit{Sofie arbejder som sygeplejerske på deltid}. Der mangler 10 ord (1--10).
  \item Skriv et ord, der passer. Du skal kun skrive \textbf{ét ord}.
\end{itemize}

\vspace{0.35em}
\hrule
\vspace{0.45em}

\textbf{\color{NCDNavy}Sofie arbejder som sygeplejerske på deltid}

\vspace{0.3em}
Sofie har arbejdet som sygeplejerske på et stort sygehus i Odense i seks år. For et år siden valgte hun at gå ned i tid, så hun nu kun arbejder 30 timer om ugen. Det betyder, at hun tjener lidt mindre end før, men til gengæld har hun fået meget mere overskud i hverdagen.

\vspace{0.25em}
På hendes afdeling er der ofte stor mangel \textbf{(1)}\ \rule{1.8cm}{0.4pt}\ uddannede sygeplejersker, og arbejdsdagen kan derfor være hektisk. Når der er travlt, bliver personalet somme tider nødt \textbf{(2)}\ \rule{1.8cm}{0.4pt}\ at løbe stærkt for at nå alle opgaverne. Tidligere følte Sofie sig ofte helt udmattet efter vagten, fordi det \textbf{(3)}\ \rule{1.8cm}{0.4pt}\ altid er let at finde ro, når man kommer sent hjem.

\vspace{0.25em}
Nu har Sofie en fast fridag om onsdagen, hvor hun kan hvile ud midt på ugen. Hun bruger dagen på at dyrke motion og hjælpe sin gamle mor, \textbf{(4)}\ \rule{1.8cm}{0.4pt}\ bor tyve minutter væk. Sammenlignet \textbf{(5)}\ \rule{1.8cm}{0.4pt}\ sine kollegaer på fuldtid har Sofie fået en langt bedre balance mellem arbejde og fritid.

\vspace{0.25em}
Efter overenskomsten har alle ansatte ret \textbf{(6)}\ \rule{1.8cm}{0.4pt}\ at søge om nedsat tid, \textbf{(7)}\ \rule{1.8cm}{0.4pt}\ det kan lade sig gøre i vagtplanen. Sofies afdelingsleder var heldigvis positiv, da hun spurgte, \textbf{(8)}\ \rule{1.8cm}{0.4pt}\ det var muligt at ændre kontrakten. I dag prioriterer hun sin fritid højere \textbf{(9)}\ \rule{1.8cm}{0.4pt}\ at undgå udbrændthed, og hun oplever langt mindre stress.

\vspace{0.25em}
Fremover håber Sofie, at flere arbejdspladser vil give mulighed for deltid. Hun er sikker på, \textbf{(10)}\ \rule{1.8cm}{0.4pt}\ fleksibel arbejdstid kan få flere sygeplejersker til at blive i faget.
\end{emailbox}

% ---------------------------------------------------------------------
\subsubsection*{Opgave 4 (6 point): 6 korte tekster（实用规定与职场消息）}

\begin{emailbox}{Kursistark --- Instruktion（Opgave 4）}
\begin{itemize}
  \item Læs de seks små tekster (1--6).
  \item Sæt \textbf{ét kryds} for det rigtige svar (A--C).
\end{itemize}
\end{emailbox}

\begin{emailbox}{1. Regler for ferie og feriepenge}
\textbf{Regler for ferie og feriepenge}\\
Som lønmodtager i Danmark optjener du 2,08 feriedage med løn eller feriepenge for hver måned, du er ansat. Det svarer til 25 feriedage på et helt ferieår, som løber fra 1. september til 31. august året efter. Du kan bruge dine optjente feriedage løbende, allerede måneden efter du har optjent dem. Hvis du har aftalt din ferie med din arbejdsgiver, skal du selv huske at bestille dine feriepenge på borger.dk, dog tidligst en måned før din første feriedag. Hvis du har været syg før din ferie begynder, har du ikke pligt til at holde ferien som planlagt, men kan få den udskudt til senere.
\par\noindent\hfill {\small\textit{Kilde: borger.dk, 2026 (red.)}}

\vspace{0.3em}
\hrule
\vspace{0.3em}
\small
\fbox{\phantom{X}}\ \textbf{A.} Du må først bestille dine feriepenge, når du er kommet tilbage fra ferie.\\[0.25em]
\fbox{\phantom{X}}\ \textbf{B.} Du kan begynde at bruge optjente feriedage allerede måneden efter, du har optjent dem.\\[0.25em]
\fbox{\phantom{X}}\ \textbf{C.} Alle lønmodtagere får automatisk deres feriepenge udbetalt den 1. september hvert år.
\end{emailbox}

\begin{emailbox}{2. A-kasse og dagpenge efter opsigelse}
\textbf{A-kasse og dagpenge efter opsigelse}\\
Hvis du bliver arbejdsløs, skal du melde dig ledig på jobnet.dk på din første ledige dag for at kunne modtage dagpenge. Du kan tidligst få udbetalt dagpenge fra den dag, du er tilmeldt jobcentret. For at have ret til fulde dagpenge som fuldtidsforsikret skal du som hovedregel have været medlem af en a-kasse i mindst ét år og have haft et vist lønindkomstniveau inden for de seneste tre år. Hvis du selv siger dit job op uden en gyldig grund, får du normalt en karantæne på tre uger, hvor du ikke modtager dagpenge.
\par\noindent\hfill {\small\textit{Kilde: star.dk, 2026 (red.)}}

\vspace{0.3em}
\hrule
\vspace{0.3em}
\small
\fbox{\phantom{X}}\ \textbf{A.} Hvis du selv siger dit arbejde op uden gyldig grund, må du vente tre uger på at få dagpenge.\\[0.25em]
\fbox{\phantom{X}}\ \textbf{B.} Du får automatisk udbetalt dagpenge, så snart du har været medlem af en a-kasse i seks måneder.\\[0.25em]
\fbox{\phantom{X}}\ \textbf{C.} Du kan først melde dig ledig på jobnet.dk, tre uger efter du er stoppet på dit arbejde.
\end{emailbox}

\begin{emailbox}{3. Kort arbejds-SMS fra kollega}
Hej Kasper\\
Vores teammøde om tirsdagen er desværre blevet flyttet fra kl. 10 til kl. 13.30, fordi Maja er til kundemøde om formiddagen. Jeg ved, du plejer at hente din datter kl. 15, så jeg har lagt vores vigtigste punkter først på dagsordenen. Skriv lige, hvis du alligevel skal gå tidligt.\\
Vh Line

\vspace{0.3em}
\hrule
\vspace{0.3em}
\small
\fbox{\phantom{X}}\ \textbf{A.} Line vil gerne aflyse tirsdagens teammøde på grund af et kundemøde.\\[0.25em]
\fbox{\phantom{X}}\ \textbf{B.} Line beder Kasper om at overtage Majas kundemøde om formiddagen.\\[0.25em]
\fbox{\phantom{X}}\ \textbf{C.} Line har tilpasset mødet, så Kasper kan nå sine vigtigste ting inden kl. 15.
\end{emailbox}

\begin{emailbox}{4. Barselsregler for nye forældre}
\textbf{Barselsregler for nye forældre}\\
Når I bliver forældre i Danmark, har I tilsammen ret til 48 ugers barselsorlov med barselsdagpenge efter barnets fødsel. Ud af disse er 11 uger øremærket til hver forælder, hvis I begge er lønmodtagere. De øremærkede uger kan ikke overføres til den anden forælder, og de skal holdes, inden barnet fylder et år --- ellers bortfalder de. De resterende uger kan I derimod frit fordele imellem jer, alt efter hvad der passer bedst til jeres familie og arbejdssituation.
\par\noindent\hfill {\small\textit{Kilde: borger.dk, 2026 (red.)}}

\vspace{0.3em}
\hrule
\vspace{0.3em}
\small
\fbox{\phantom{X}}\ \textbf{A.} Begge forældre kan frit give alle deres barselsuger videre til deres partner.\\[0.25em]
\fbox{\phantom{X}}\ \textbf{B.} De 11 øremærkede uger går tabt, hvis man ikke bruger dem i barnets første leveår.\\[0.25em]
\fbox{\phantom{X}}\ \textbf{C.} Hver forælder har ret til 48 ugers barselsorlov med barselsdagpenge alene.
\end{emailbox}

\begin{emailbox}{5. Kort SMS om feedback på projektrapport}
Hej Nadia\\
Jeg har nu læst dit udkast til projektrapporten igennem. Indledningen og tallene ser rigtig stærke ud, men jeg synes, at konklusionen godt kunne være lidt skarpere. Kan vi kigge på den sammen over en kop kaffe i morgen kl. 9, før vi sender rapporten til kunden?\\
Mvh Thomas

\vspace{0.3em}
\hrule
\vspace{0.3em}
\small
\fbox{\phantom{X}}\ \textbf{A.} Thomas mener, at hele rapporten allerede er klar til at blive sendt til kunden.\\[0.25em]
\fbox{\phantom{X}}\ \textbf{B.} Thomas foreslår, at Nadia helt skriver en ny indledning til projektrapporten.\\[0.25em]
\fbox{\phantom{X}}\ \textbf{C.} Thomas vil gerne gennemgå den sidste del af rapporten med Nadia i morgen.
\end{emailbox}

\begin{emailbox}{6. Regler for lån og gebyrer på folkebiblioteket}
\textbf{Regler for lån og gebyrer på folkebiblioteket}\\
Når du låner bøger, film eller brætspil på folkebiblioteket, er lånetiden normalt 30 dage. Hvis ingen andre har reserveret materialet, kan du forny dit lån op til tre gange via bibliotekets app eller hjemmeside. Du får automatisk en påmindelse på sms eller e-mail tre dage før afleveringsfristen. Afleverer du ikke materialet til tiden, skal voksne over 18 år betale et gebyr, som stiger, jo flere dage fristen overskrides. Børn og unge under 18 år betaler derimod ikke gebyr ved for sen aflevering, men får en erstatningsregning, hvis materialet bliver væk.
\par\noindent\hfill {\small\textit{Kilde: bibliotek.dk, 2026 (red.)}}

\vspace{0.3em}
\hrule
\vspace{0.3em}
\small
\fbox{\phantom{X}}\ \textbf{A.} Børn under 18 år skal kun betale penge, hvis en lånt ting bliver helt væk.\\[0.25em]
\fbox{\phantom{X}}\ \textbf{B.} Man kan altid forny alle lånte bøger tre gange i bibliotekets app.\\[0.25em]
\fbox{\phantom{X}}\ \textbf{C.} På folkebiblioteket skal både børn og voksne betale et gebyr, hvis de afleverer for sent.
\end{emailbox}

% =====================================================================
% DEL 2: SKRIVNING (Dual Tasks)
% =====================================================================
\ncdsubsec{Del 2: Skrivning（书面表达｜Opgave 1 私人回信 \& Opgave 2 论坛辩论贴）}

\subsubsection*{Opgave 1（私人邮件回复 · Minimum 125 ord）}

\textbf{Situation:} Du har fået en e-mail fra din ven, Mikkel.\\
\textbf{Opgave:} Læs e-mailen, og skriv et svar til Mikkel. Du skal svare på Mikkels spørgsmål.\\
\textbf{Længde:} Minimum 125 ord.

\begin{emailbox}{Modtaget e-mail fra Mikkel}
Hej.

\vspace{0.3em}
Hvordan går det hos jer? Jeg har tænkt på dig, for sidst vi snakkede, havde du jo lige travlt med at pakke flyttekasser. Hvordan er det gået med at flytte ind i den nye lejlighed? Er I faldet godt til i det nye kvarter? Jeg er spændt på at høre, hvordan det hele fungerer.

\vspace{0.3em}
Her hos mig er der sket noget spændende på arbejdet. Jeg er lige blevet tilbudt en stilling som teamleder i min afdeling! Som en del af aftalen kan jeg vælge mellem to goder: enten en højere månedsløn eller en ekstra uges ferie hvert år. Jeg er virkelig i tvivl om, hvad der bedst kan betale sig. Hvad synes du, jeg skal vælge?

\vspace{0.3em}
Der er også en anden ting, jeg går og tænker over. På min nye arbejdsplads er det ret normalt, at både ledere og kollegaer svarer på arbejdsmails i weekenden. Det har fået mig til at overveje, om man bør acceptere det, eller om man helt skal holde fri i fritiden. Hvad mener du om den kultur?

\vspace{0.3em}
\ldots
\end{emailbox}

\begin{emailbox}{Dit svar｜Til: Mikkel \qquad Emne: Sv: Flytning, teamleder og arbejdskultur}
\vspace{0.3em}
\rule{\linewidth}{0.4pt}\\[0.55em]
\rule{\linewidth}{0.4pt}\\[0.55em]
\rule{\linewidth}{0.4pt}\\[0.55em]
\rule{\linewidth}{0.4pt}\\[0.55em]
\rule{\linewidth}{0.4pt}\\[0.55em]
\rule{\linewidth}{0.4pt}\\[0.55em]
\rule{\linewidth}{0.4pt}\\[0.55em]
\rule{\linewidth}{0.4pt}\\[0.55em]
\rule{\linewidth}{0.4pt}\\[0.55em]
\rule{\linewidth}{0.4pt}\\[0.55em]
\rule{\linewidth}{0.4pt}\\[0.55em]
\rule{\linewidth}{0.4pt}
\end{emailbox}

\subsubsection*{Opgave 2（观点辩论帖 / \textit{Blogindlæg} · Minimum 100 ord）}

\textbf{Situation:} Henrik Lind har skrevet et blogindlæg på en hjemmeside, hvor man diskuterer aktuelle emner.\\
\textbf{Opgave:} Skriv et blogindlæg, hvor du diskuterer Henrik Linds synspunkt om, at hjemmearbejde ødelægger fællesskabet på arbejdspladsen. Er du enig? Hvorfor? / Hvorfor ikke?\\
\textbf{Længde:} Minimum 100 ord.

\begin{emailbox}{Blogindlæg af Henrik Lind：\textit{Hjemmearbejde ødelægger fællesskabet}}
\textbf{\color{NCDNavy}Hjemmearbejde ødelægger fællesskabet}

\vspace{0.3em}
Siden corona har det været helt almindeligt, at folk arbejder hjemmefra flere dage om ugen. Selvom det lyder praktisk for den enkelte, oplever jeg på min egen arbejdsplads, at hjemmearbejdet langsomt ødelægger fællesskabet. Kontoret er ofte halvtomt, og de små snakke ved kaffemaskinen er næsten forsvundet. Nye kollegaer har svært ved at lære kulturen at kende og få hurtig hjælp, når halvdelen af teamet kun sidder bag en skærm derhjemme. Et godt arbejdsmiljø kræver, at vi mødes fysisk og skaber noget sammen. Derfor mener jeg, at virksomheder bør sætte klare grænser og kræve, at medarbejderne som udgangspunkt møder ind på arbejdspladsen hver dag.
\end{emailbox}

\begin{emailbox}{Dit blogindlæg｜Diskussion af Henrik Linds synspunkt}
\vspace{0.3em}
\rule{\linewidth}{0.4pt}\\[0.55em]
\rule{\linewidth}{0.4pt}\\[0.55em]
\rule{\linewidth}{0.4pt}\\[0.55em]
\rule{\linewidth}{0.4pt}\\[0.55em]
\rule{\linewidth}{0.4pt}\\[0.55em]
\rule{\linewidth}{0.4pt}\\[0.55em]
\rule{\linewidth}{0.4pt}\\[0.55em]
\rule{\linewidth}{0.4pt}\\[0.55em]
\rule{\linewidth}{0.4pt}\\[0.55em]
\rule{\linewidth}{0.4pt}\\[0.55em]
\rule{\linewidth}{0.4pt}
\end{emailbox}

% =====================================================================
% DEL 3: MUNDTLIG KOMMUNIKATION
% =====================================================================
\ncdsubsec{Del 3: Mundtlig kommunikation（口语交际全真题卡｜Opgave 1 \& 2）}

\subsubsection*{Mundtlig Opgave 1 (Emneopgave｜Del 1: 1--2 min. + Del 2: 3--4 min.)}

\begin{tipbox}{DU3 Modul 4 官方口试规则（SIRI / Modultest.dk \textit{Opgave 1: Emneopgave}）}
\small
\begin{itemize}
  \item \textbf{Del 1}：\textit{Du vælger ved lodtrækning ét af de to emner, som du har forberedt. Herefter fortæller du om emnet i 1--2 minutter.}（考生须事先围绕 \textbf{danske forhold} 准备两个主题，考场抽签决定其一并脱稿连贯陈述 1--2 分钟。）
  \item \textbf{Del 2}：\textit{Præsentationen følges op af spørgsmål til emnet i 3--4 minutter. Tester stiller spørgsmål, der lægger op til, at du forklarer, uddyber og begrunder i lidt længere taleture, fx: Vil du fortælle lidt mere om …? / Kan du give nogen eksempler på …? / Hvad er fordelene / ulemperne ved …? Hvorfor? / Hvad synes du om …? Hvorfor?}
  \item \textbf{纯关键词（stikord）铁律}：官方明确规定只能携带只含单词/短语的提示卡（\textit{stikord, dvs. enkeltord og ikke hele sætninger}），严禁写完整句子照读；且内容必须同时兼顾\textbf{宏观丹麦社会现状（generelt fokus）}与\textbf{个人亲身体验（egne erfaringer）}，整体评估为 \textit{Bestået / Ikke bestået}。
\end{itemize}
\end{tipbox}

\begin{grammatikbox}{Stikord-kort 1A（已填纯关键词卡）：Arbejdsliv: balance mellem job og fritid}
\small
\renewcommand{\arraystretch}{1.3}
\begin{tabularx}{\linewidth}{@{} >{\leavevmode\bfseries\color{NCDNavy}\raggedright\arraybackslash}p{4.3cm} >{\raggedright\arraybackslash}X @{}}
\toprule
\multicolumn{2}{c}{\normalsize\bfseries\color{NCDRed} Emne 1A: Arbejdsliv med fokus på balance mellem job og fritid \quad\small\color{NCDGray}（仅含 stikord）} \\
\midrule
1. Generelt fokus（丹麦职场常态） & 37 timer/uge · fladt hierarki · tillid frem for kontrol · hente børn kl. 16 \\
\midrule
2. Fordele ved fleksibilitet & mere overskud · mindre stress · tid til motion og familie · færre sygedage \\
\midrule
3. Ulemper / udfordringer & uklar grænse job/privatliv · arbejdsmail på telefonen · forventning om at tjekke \\
\midrule
4. Egne erfaringer \& vaner & arbejde effektivt i dagtimerne · slukke mail kl. 17 · løbetur \& aftensmad med familien \\
\midrule
5. Fremtidsperspektiv & ønske om 4 dages arbejdsuge · tre fridage til hvile · højere arbejdsglæde \\
\bottomrule
\end{tabularx}

\vspace{0.45em}
\textbf{\color{NCDRed}Del 2 考官追问预测（Eksaminators opfølgende spørgsmål · 4 题）：}
\begin{enumerate}
  \item \textit{Vil du fortælle lidt mere om, hvordan du selv oplever det flade hierarki og tilliden på en dansk arbejdsplads?}
  \item \textit{Kan du give nogle eksempler på, hvordan du sætter grænser mellem arbejde og fritid i din egen hverdag?}
  \item \textit{Hvad er efter din mening fordelene og ulemperne ved at arbejde hjemmefra et par dage om ugen?}
  \item \textit{Hvad synes du om forskellen mellem arbejdskulturen i Danmark og i dit hjemland? Hvorfor?}
\end{enumerate}
\end{grammatikbox}

\begin{grammatikbox}{Stikord-kort 1B（已填纯关键词卡）：Ligestilling: barsel og karriere}
\small
\renewcommand{\arraystretch}{1.3}
\begin{tabularx}{\linewidth}{@{} >{\leavevmode\bfseries\color{NCDNavy}\raggedright\arraybackslash}p{4.3cm} >{\raggedright\arraybackslash}X @{}}
\toprule
\multicolumn{2}{c}{\normalsize\bfseries\color{NCDRed} Emne 1B: Ligestilling med fokus på barsel og karriere \quad\small\color{NCDGray}（仅含 stikord）} \\
\midrule
1. Danske regler（制度背景） & 48 ugers barsel tilsammen · barselsdagpenge · 11 øremærkede uger pr. forælder (2022) \\
\midrule
2. Formål \& fordele & ligestilling hjemme og på jobbet · tættere bånd far--barn · ikke kun kvinders karrierepause \\
\midrule
3. Ulemper / bekymringer & frygt for løntab eller færre forfremmelser · tilpasning i små virksomheder · mindre valgfrihed \\
\midrule
4. Egne observationer & danske fædre med barnevogn · stor forskel fra hjemlandet · længere orlov for mænd \\
\midrule
5. Vurdering（个人立场） & vigtigt skridt mod reel ligestilling · ændrer holdninger på arbejdsmarkedet \\
\bottomrule
\end{tabularx}

\vspace{0.45em}
\textbf{\color{NCDRed}Del 2 考官追问预测（Eksaminators opfølgende spørgsmål · 4 题）：}
\begin{enumerate}
  \item \textit{Vil du fortælle lidt mere om, hvorfor du mener, at øremærket barsel til fædre er vigtig?}
  \item \textit{Kan du give nogle eksempler på, hvordan arbejdsdelingen i familien kan påvirke kvinders og mænds karriere?}
  \item \textit{Hvad er fordelene og ulemperne ved, at staten bestemmer, at 11 uger ikke kan overføres til partneren?}
  \item \textit{Hvad synes du om barselsreglerne i Danmark sammenlignet med reglerne i dit hjemland? Hvorfor?}
\end{enumerate}
\end{grammatikbox}

% ---------------------------------------------------------------------
\subsubsection*{Mundtlig Opgave 2 (Del 1 Individuel ca. 3 min. pr. kursist + Del 2 Paropgave ca. 4 min.)}

\begin{grammatikbox}{Opgaveark A \& B --- EMNE: Arbejde og fritid（个人图片情景提示卡）}
\small
\begin{tabularx}{\linewidth}{@{} X @{\hspace{0.4cm}} X @{}}
\begin{tcolorbox}[enhanced, colback=white, colframe=NCDTeal, boxrule=0.8pt, arc=3pt, left=7pt, right=7pt, top=6pt, bottom=6pt]
  \textbf{\color{NCDNavy}Opgaveark A：居家远程办公}\\
  {\footnotesize\color{NCDGray}[图片 A：一名职场人士坐在明亮的自家客厅餐桌或书桌前，面前放着打开的笔记本电脑与咖啡杯，窗外是宁静的住宅区。]}
\end{tcolorbox}
&
\begin{tcolorbox}[enhanced, colback=white, colframe=NCDTeal, boxrule=0.8pt, arc=3pt, left=7pt, right=7pt, top=6pt, bottom=6pt]
  \textbf{\color{NCDNavy}Opgaveark B：兼职父母提前接娃}\\
  {\footnotesize\color{NCDGray}[图片 B：一位兼职工作的年轻父母下午三点左右就来到幼儿园门口接孩子，牵着孩子的手开心地走在回家的路上。]}
\end{tcolorbox}
\end{tabularx}
\end{grammatikbox}

\begin{emailbox}{TESTERARK --- 官方考官引导语、必问三题与双人辩论题卡（EMNE: Arbejde og fritid）}
\small\textcolor{NCDGray}{\textit{Instruktionen skal følges af tester. Hvis der vælges en anden sproglig formulering, skal instruktioner og spørgsmål indholdsmæssigt og sprogligt ligge så tæt på forlægget som muligt.}}

\vspace{0.35em}
\textbf{\color{NCDRed}Del 1 (individuel opgave -- ca. 3 min. pr. kursist)}\\
\textbf{\color{NCDNavy}KURSIST A}\\
Opgaveark A lægges foran kursist A, der præsenteres for emnet på følgende måde:
\begin{quote}
\textit{Emnet er 'Arbejde og fritid'. Her er et billede, som viser en person, der arbejder hjemmefra.}
\end{quote}
Tester stiller derpå nedenstående obligatoriske spørgsmål:
\begin{enumerate}
  \item \textit{Der er nogle danskere, der arbejder hjemmefra flere dage om ugen. Hvorfor tror du, de gør det?}
  \item \textit{Hvilke fordele kan der være ved det?}
  \item \textit{Mener du, der kan være nogen ulemper ved at arbejde hjemmefra? Hvorfor? / Hvorfor ikke?}
\end{enumerate}

\textbf{\color{NCDNavy}KURSIST B}\\
Opgaveark B lægges foran kursist B, der præsenteres for emnet på følgende måde:
\begin{quote}
\textit{Emnet er 'Arbejde og fritid'. Her er et billede, som viser en ung forælder på deltid, der henter sit barn tidligt fra børnehave.}
\end{quote}
Tester stiller derpå nedenstående obligatoriske spørgsmål:
\begin{enumerate}
  \item \textit{Der er nogle forældre, der vælger at arbejde på deltid, mens deres børn er små. Hvorfor tror du, de gør det?}
  \item \textit{Hvilke fordele kan der være ved det?}
  \item \textit{Mener du, der kan være nogen ulemper ved at arbejde på deltid? Hvorfor? / Hvorfor ikke?}
\end{enumerate}

\vspace{0.35em}
\hrule
\vspace{0.35em}
\textbf{\color{NCDRed}Del 2 (paropgave -- ca. 4 min.)}\\
Testen fortsætter som en samtale mellem de to kursister, hvor de udveksler synspunkter om et emne og begrunder disse. Tester præsenterer emnet med følgende input til samtalen:
\begin{quote}
\textit{Flere virksomheder i Danmark har indført en fire dages arbejdsuge med fuld løn. Hvad synes I om det? I skal nu diskutere, hvilke fordele og ulemper, I mener, der kan være for den enkelte og for samfundet. I skal stille spørgsmål til hinanden og svare på hinandens spørgsmål.}
\end{quote}
\small\textcolor{NCDGray}{Tester fungerer som mediator og bryder ind, hvis samtalen ikke handler om emnet, eller hvis samtalen går i stå eller overtages helt af den ene kursist.}
\end{emailbox}

% =====================================================================
% DEL 4: RETTEARK, DETALJERET ANALYSE & HØJSCORE-MODELLER
% =====================================================================
\clearpage
\ncdsubsec{Del 4: Retteark, Detaljeret Analyse \& Højscore-Modeller}

\subsubsection*{4.1 官方 \textit{Retteark} 评分卡（29 point）与考前自查准则}

\small
\renewcommand{\arraystretch}{1.26}
\begin{longtable}{@{\extracolsep{\fill}} >{\leavevmode\bfseries\color{NCDNavy}\raggedright\arraybackslash}p{3.1cm} >{\raggedright\arraybackslash}p{8.6cm} >{\centering\arraybackslash}p{1.6cm} >{\centering\arraybackslash}p{2.2cm} @{}}
\toprule
\textbf{考核题型} & \textbf{Rettenøgle（标准答案与官方可接受变体）} & \textbf{分值} & \textbf{自测 / 评估} \\
\midrule
\endfirsthead
\toprule
\textbf{考核题型（续表）} & \textbf{Rettenøgle（标准答案与官方可接受变体）} & \textbf{分值} & \textbf{自测 / 评估} \\
\midrule
\endhead
\bottomrule
\endfoot
\textbf{Læsning Opgave 1}\newline（8 题） &
1. \textbf{B} (\textit{gevinster}) \quad 2. \textbf{A} (\textit{udkørte}) \quad 3. \textbf{D} (\textit{indføre}) \quad 4. \textbf{C} (\textit{Derimod})\newline
5. \textbf{A} (\textit{udfordring}) \quad 6. \textbf{D} (\textit{nødt til}) \quad 7. \textbf{B} (\textit{Derudover}) \quad 8. \textbf{C} (\textit{Alt i alt})\newline
\textcolor{NCDGray}{\small （例题 0 为 C \textit{fristende}）} & 8 point & \rule{0.65cm}{0.4pt} / 8 \\
\midrule
\textbf{Læsning Opgave 2}\newline（5 题） &
1. \textbf{E} \quad 2. \textbf{G} \quad 3. \textbf{C} \quad 4. \textbf{H} \quad 5. \textbf{D}\newline
\textcolor{NCDGray}{\small （例句 0 为 A｜多余干扰项：B, F）} & 5 point & \rule{0.65cm}{0.4pt} / 5 \\
\midrule
\textbf{Læsning Opgave 3}\newline（10 题 · 无词库） &
1. \textbf{på} \quad 2. \textbf{til} \quad 3. \textbf{ikke} \quad 4. \textbf{som} eller \textbf{der} \quad 5. \textbf{med}\newline
6. \textbf{til} \quad 7. \textbf{hvis} (eller \textbf{såfremt} / \textbf{når}) \quad 8. \textbf{om} (eller \textbf{hvorvidt})\newline
9. \textbf{for} \quad 10. \textbf{at} & 10 point & \rule{0.65cm}{0.4pt} / 10 \\
\midrule
\textbf{Læsning Opgave 4}\newline（6 题） &
1. \textbf{B} \quad 2. \textbf{A} \quad 3. \textbf{C} \quad 4. \textbf{B} \quad 5. \textbf{C} \quad 6. \textbf{A} & 6 point & \rule{0.65cm}{0.4pt} / 6 \\
\midrule
\textbf{\color{NCDRed}Læsning 总计} & \textbf{Opgave 1--4 合订总评（50 分钟 · 禁用词典）} & \textbf{29 point} & \rule{0.65cm}{0.4pt} / 29 \\
\midrule
\textbf{Skrivning Opgave 1}\newline（私人邮件回信） &
词数严格 \textbf{$\ge$ 125 ord}；全面回应 Mikkel 全部 3 个维度（搬家近况、加薪 vs 多休一周假选择与理由、周末回复工作邮件文化观点）；自然融入情态小品词。 & 整体评估 & Bestået / Ikke bestået \\
\midrule
\textbf{Skrivning Opgave 2}\newline（观点辩论帖） &
词数严格 \textbf{$\ge$ 100 ord}；清晰表明是否同意 Henrik Lind 的观点，使用让步反驳并给出严谨理由与折中方案。 & 整体评估 & Bestået / Ikke bestået \\
\midrule
\textbf{Mundtlig Opgave 1}\newline（Emneopgave） &
仅带纯 \textit{stikord} 提示卡；Del 1 陈述 1--2 分钟兼顾丹麦社会与个人经验；Del 2 以较长话轮回答追问并明确论证利弊与原因。 & 整体评估 & Bestået / Ikke bestået \\
\midrule
\textbf{Mundtlig Opgave 2}\newline（Ind. + Paropgave） &
Del 1（约 3 分钟/人）：回答图片三问（原因 / 优点 / 缺点与理由）；Del 2（约 4 分钟）：双人互动互问互答，必须兼顾\textbf{对个人与对社会}两方面利弊。 & 整体评估 & Bestået / Ikke bestået \\
\end{longtable}
\normalsize

\begin{tipbox}{考前自查准则与失分诊断（Self-Check Checklist）}
\small
\begin{itemize}
  \item \textbf{总分达标线参考}：官方 \textit{Retteark} 仅记 29 分原始总分，各语言中心依 SIRI 指引划定的及格线通常约为 60--65\%（即 \textbf{18--19 分以上}）。建议冲刺自测稳定拿到 \textbf{20+/29}。
  \item \textbf{Opgave 3 无词库开放完形失分排查}：若在 Opgave 3 低于 7/10，请重点检查三类失分：\ding{172} 固定介词搭配（\textit{mangel på, nødt til, sammenlignet med, ret til}）；\ding{173} 从句连接词精准区分（\textit{om} 间接是非问 vs \textit{hvis} 条件 vs \textit{at} 宾语从句）；\ding{174} 否定词 \textit{ikke} 在从句“主语后、动词前”的定位（如 \textit{fordi det ikke altid er...}）。
  \item \textbf{写作与口语自查红线}：\ding{172} 写完务必逐词点数，确保 Opgave 1 $\ge$ 125 词、Opgave 2 $\ge$ 100 词；\ding{173} 口语 \textit{Emneopgave} 严禁带完整句子进考场；\ding{174} 口语 Opgave 2 Del 2 必须主动发问并同时讨论 \textit{for den enkelte}（对个人）与 \textit{for samfundet}（对社会）。
\end{itemize}
\end{tipbox}

% ---------------------------------------------------------------------
\subsubsection*{4.2 阅读 Opgave 1--4 逐题解析、全文中文精译与高频词表}

\begin{answerbox}
\textbf{\color{NCDNavy}【一、Opgave 1（Fire dages arbejdsuge i praksis）逐题语法与语义解析】}
\begin{itemize}
  \item \textbf{(0) C (\textit{fristende}) [Eksempel]}：后文 \textit{fordi medarbejderne får en ekstra fridag uden at gå ned i løn}（不降薪多休一天）说明这个想法听起来“诱人”（\textit{fristende}）。
  \item \textbf{(1) B (\textit{gevinster})}：后文列举员工更专注、少开会、精力更足、病假下降，均为积极“收益/好处”（\textit{gevinster ved}）。\textbf{排雷}：A（缺点）、C（成本）、D（限制）均为负面词。
  \item \textbf{(2) A (\textit{udkørte})}：与 \textit{møder ind med mere energi}（带着更多精力来上班）并列，表示周末后“不那么疲惫/精疲力竭”（\textit{mindre udkørte}）。B（自豪）、C（投入）、D（有准备）均与休假恢复精力逻辑不符。
  \item \textbf{(3) D (\textit{indføre})}：搭配 \textit{en fire dages arbejdsuge ... på alle arbejdspladser}，表示“在所有工作场所推行/实施四天工作制”。A（取消）、B（放弃）、C（禁止）与第三段“办公室易于实行，医院/幼儿园难实行”对比矛盾。
  \item \textbf{(4) C (\textit{Derimod})}：前句讲办公室/IT 企业很容易在四天内完成同等工作，本句讲医院、幼儿园则困难得多，为典型对比转折“相反/另一方面”（\textit{Derimod er det langt sværere...}），且位于句首触发 V2 倒装。\textbf{排雷}：\textit{Nemlig} 不能放在句首占位；\textit{Derfor / Desuden} 逻辑方向错误。
  \item \textbf{(5) A (\textit{udfordring})}：后文紧接着讲工作压力、忙碌和焦虑，故此处是“另一个挑战”（\textit{En anden udfordring}），呼应前一段“在所有行业推行并不容易”的第一个难点。
  \item \textbf{(6) D (\textit{nødt til})}：固定搭配 \textit{være nødt til at gøre noget}（不得不做某事），表示如果不减量，员工不得不加班或下班后回邮件才能做完。A（害怕）、B（高兴）、C（自豪）均与 \textit{for at nå det hele}（为了赶完所有工作）背离。
  \item \textbf{(7) B (\textit{Derudover})}：第四段讲工作压力，第五段补充劳工市场研究者的另一视角“不同行业的灵活性各不相同”，故用递进副词 \textit{Derudover}（此外）。
  \item \textbf{(8) C (\textit{Alt i alt})}：全篇收束句，总括“总体而言，只要任务量调整得切合实际，四天工作制有望成为越来越多现代企业的常规选项”。A（例如）、B（相反）、D（遗憾的是）均不适合总括积极前景。
\end{itemize}
\end{answerbox}

\begin{transbox}
\small
\textbf{\color{NCDTeal}[Opgave 1 全文中文精译]}\\
大多数丹麦人每周依然工作约 37 小时、分布于五个工作日，但近年来已有多家企业和少数市政部门试行每周仅工作四天的制度。这个主意听起来或许很诱人，因为员工可以在不降薪的情况下多获得一天休息日。\\[3pt]
试行经验确实显示出缩短工作周所带来的诸多明显收益。当员工知道周五放假时，许多人会更擅长优先处理最重要的任务并缩短会议时间。多位管理者反映，员工周一上班时精力更充沛，周末过后也不那么疲惫。同时，病假缺勤率往往会下降，因为人们有了更多时间陪伴家人和充分休息。\\[3pt]
尽管如此，四天工作制并非在所有工作场所都同样容易推行。在办公室或 IT 企业里，人们往往能更高效地工作，在四天内完成同样的工作量。相反，在医院、幼儿园或居家护理机构中则困难得多，因为这些岗位必须始终有人在场照顾市民。\\[3pt]
另一个挑战则是那四个工作日里的工作压力。如果企业不减少任务总量，这四天就可能迅速变得过于忙碌且令人倍感压力。部分员工感到，自己不得不晚上加班或在业余时间回复邮件，才能把所有事情做完。\\[3pt]
此外，多位劳动力市场研究人员指出，灵活性在不同行业中表现并不相同。对某些人来说三天小长假最合适，而另一些人则更希望缩短每天工时或拥有居家办公的选择。总体而言，只要工作任务能够切合实际地调整，四天工作制有望在越来越多现代企业中成为一项固定制度。
\end{transbox}

\begin{answerbox}
\textbf{\color{NCDNavy}【二、Opgave 2（Fra bankrådgiver til snedker）逐题衔接解析】}
\begin{itemize}
  \item \textbf{(0) A [Eksempel]}：承接前句的“隐隐不安（\textit{en stille uro}）”，具体交代他几乎整天坐在屏幕前面对表格、会议和数字，顺理成章引出下句“每晚合上笔记本电脑时总觉得一天下来没创造任何具体成果”。
  \item \textbf{(1) E}：前句说他“本来只想学修一张旧餐桌（\textit{Egentlig ville han bare...}）”，\textbf{E} 用转折连词 \textit{Men} 接住：“但工坊里的体验在他心中点燃了意想不到的热情”，随即自然接出下句“新刨橡木的独特香气与亲手塑形木材的喜悦立刻击中了他”。
  \item \textbf{(2) G}：前句交代他辞去银行铁饭碗去读木工职业教育，\textbf{G} 精准概括“从固定银行工资过渡到学徒工资，然而需要大幅调整个人财务（\textit{Overgangen fra fast bankløn til elevløn...}）”，直接解释后句的“作为成年学徒收入锐减，只好重编预算、停掉昂贵出国游”。
  \item \textbf{(3) C}：前句描述日常既有精细活也有重体力搬运，\textbf{C} 紧接身体感受“一开始每天下午他的背和手都酸痛”，接着与后句的 \textit{Samtidig}（同时，工坊师傅却很欣赏他原本的财务与客户沟通能力）形成“体力吃力 vs 经验加分”的精妙呼应。
  \item \textbf{(4) H}：前句介绍他开了工坊主打可持续丹麦木材耐用家具，\textbf{H} 指出“尽管大型连锁家具店竞争激烈，这家小公司订单已经应接不暇”，而下一句立刻用 \textit{nemlig}（因为……嘛）解释原因：“顾客愿意多花一点钱直接与手艺人沟通并了解材料来源”。
  \item \textbf{(5) D}：前句总结“成功不仅关乎工资单，而是周一早上期待去上班”，\textbf{D} 以 \textit{Derfor} 导出“因此他常鼓励感到在工作里陷入僵局的人考虑换赛道”，完美引出最后的直接引语。
  \item \textbf{多余干扰项排雷}：\textbf{B} 称“依然怀念金融界的高额奖金与快节奏”，与末段“丝毫不后悔”完全矛盾；\textbf{F} 谈银行向农村买房年轻人提供低息贷款，纯属利用“银行”“乡下（Odder 附近）”关键词制造的离题陷阱。
\end{itemize}
\end{answerbox}

\begin{transbox}
\small
\textbf{\color{NCDTeal}[Opgave 2 全文中文精译]}\\
安德斯·霍尔姆快三十岁时，他的未来看上去早已完全定型。他完成过扎实的经济学教育，并在奥胡斯一家银行担任理财顾问五年。薪水不错、同事也很友善，但每当望向窗外的步行街，他内心总有一丝静静滋长的不安。\textbf{（0-A）他几乎整天都在屏幕前处理电子表格、会议和数字。}每天晚上合上笔记本电脑时，他都觉得自己这一天根本没有创造出任何具体实在的东西。\\[3pt]
转折点发生在十一月一个下雨的周末，当时安德斯在当地夜校报了一个短期家具木工晚间班。其实他本来只想学着修好自己公寓里的一张旧餐桌。\textbf{（1-E）然而在工坊里的那段体验却在他内心点燃了一股意想不到的热情。}新刨出的橡木散发出的独特木香，以及用双手塑造木头的快乐瞬间击中了他。才上了三个晚上，他就已经开始从图书馆借阅关于木种、榫卯拼合与工具的专业书籍了。\\[3pt]
半年后，安德斯作出了一个让家人和朋友都大吃一惊的勇敢决定。他辞去了银行里安稳的职位，开始攻读木匠职业教育。\textbf{（2-G）然而从固定银行薪水过渡到学徒工资，需要在个人财务上作出相当大的调整。}作为成年学徒，他的收入一下子比以前少了很多，因此他和女友不得不彻底重做预算，并在好几年里放弃了昂贵的出国旅行。\\[3pt]
幸运的是，安德斯很快在奥泽（Odder）外的一家小型本地工坊找到了学徒岗位。这里的日常既有精细的精密加工，也有搬运重物，需要为私人客户定制厨房与书架。\textbf{（3-C）起初每天下午他的后背和双手都酸痛不已，因为身体必须逐步适应体力劳动。}与此同时，工坊的师傅也非常欣喜能有一位熟悉财务、报价单制作及客户沟通的学徒，这些都得益于他过往的职场经历。\\[3pt]
如今，在经历那次重大的职业转型四年后，安德斯已是正式出师的家具木匠，并与一位昔日同学合开了自己的工坊。他们专注于使用来自可持续森林的丹麦木材制作耐用的家具。\textbf{（4-H）尽管来自大型家具连锁店的竞争十分激烈，这家小公司却已经接满了订单。}原来，只要能直接与手艺人交谈并确切看到原料的来源，顾客们都愿意多支付一点费用。\\[3pt]
每当安德斯回望自己在银行的岁月，他对当初的学历与这次巨大改变都毫无悔意。这次转行让他明白，成功不仅关乎工资条和头衔地位，而在于周一早晨会满怀期待地去上班。\textbf{（5-D）因此，如果别人觉得自己在工作上陷入了僵局，他经常鼓励对方认真考虑一次职业转轨。}“最重要的是，在日常生活变成没有方向的惯性之前，人要敢于倾听自己内心的声音，”他微笑着说道。
\end{transbox}

\begin{answerbox}
\textbf{\color{NCDNavy}【三、Opgave 3（Sofie arbejder som sygeplejerske på deltid）无词库填词解析】}
\begin{itemize}
  \item \textbf{(1) på}：固定名词 + 介词搭配 \textit{mangel på noget}（缺乏某物）。
  \item \textbf{(2) til}：固定搭配 \textit{blive nødt til at + 不定式}（不得不……）。
  \item \textbf{(3) ikke}：原因从句 \textit{fordi det ikke altid er let...} 考查否定词 \textit{ikke} 的精确位置（主语 \textit{det} 之后、变位动词 \textit{er} 之前，搭配 \textit{ikke altid}，此处 \textit{aldrig / ingen} 均不合语法）。
  \item \textbf{(4) som / der}：定语从句关系代词作主语，修饰 \textit{sin gamle mor}，因作主语故 \textit{som} 与 \textit{der} 均完全正确。
  \item \textbf{(5) med}：固定过去分词介词短语 \textit{sammenlignet med}（与……相比）。
  \item \textbf{(6) til}：固定搭配 \textit{have ret til at + 不定式}（有权利做某事）。
  \item \textbf{(7) hvis / såfremt / når}：条件从句连词，引出“如果/只要排班表里允许的话”。
  \item \textbf{(8) om / hvorvidt}：动词 \textit{spurgte} 后接间接是非问句，必须用 \textit{om}（是否），切忌误写 \textit{hvis}！正式书面亦可填 \textit{hvorvidt}。
  \item \textbf{(9) for}：目的不定式结构 \textit{for at undgå udbrændthed}（为了避免职业倦怠）。
  \item \textbf{(10) at}：固定短语 \textit{være sikker på, at...} 引导宾语从句，不可漏填。
\end{itemize}
\end{answerbox}

\begin{transbox}
\small
\textbf{\color{NCDTeal}[Opgave 3 全文中文精译]}\\
索菲已在欧登塞的一家大型医院当了六年护士。一年前，她选择减少工时，因此现在每周只工作 30 小时。这意味着她挣得比以前略少一些，但作为回报，她在日常生活中收获了充沛得多的精力与从容。\\[3pt]
在她所在的科室里，受过专业培训的护士经常严重短缺，因此工作日可能会相当紧张忙碌。工作繁忙时，医护人员有时不得不快步奔走以完成所有任务。过去，索菲在轮班后经常感到精疲力竭，因为晚归到家后并不总是那么容易平静下来。\\[3pt]
如今索菲每周三都有一个固定的休息日，可以在周中好好休整一番。这一天她用来锻炼身体，并照顾住在二十分钟路程外的年迈母亲。与全职同事相比，索菲在工作与休闲之间找到了好得多的平衡。\\[3pt]
根据劳资集体协议，只要排班表允许，所有员工都有权申请缩短工时。当索菲询问是否可以修改合同时，她的科室主管幸运地给予了积极回应。如今她更优先考虑自己的业余时间，以避免职业倦怠，感到的压力也少得多。\\[3pt]
展望未来，索菲希望更多工作场所能提供兼职机会。她确信，更灵活的工作时间能让更多经验丰富的护士长期留在这个行业。
\end{transbox}

\begin{answerbox}
\textbf{\color{NCDNavy}【四、Opgave 4（6 则实用文本）定位与排雷解析】}
\begin{itemize}
  \item \textbf{1. B}：原文明确写 \textit{Du kan bruge dine optjente feriedage løbende, allerede måneden efter du har optjent dem} 与 B 完全同义。\textbf{排雷}：A 错把 \textit{tidligst en måned før}（最早休假前一个月）偷换为“回来后才能申请”；C 错称“9 月 1 日自动发放”，原文要求 \textit{selv huske at bestille}（自行申请）。
  \item \textbf{2. A}：原文末句 \textit{Hvis du selv siger dit job op uden en gyldig grund, får du normalt en karantæne på tre uger, hvor du ikke modtager dagpenge} 与 A 精准对应。\textbf{排雷}：B 把“至少一年（\textit{mindst ét år}）”改错为“六个月”且忽略登记要求；C 把“三周等待期”张冠李戴为“三周后才能在 jobnet 登记”。
  \item \textbf{3. C}：Line 把重要议题排在议程最前面（\textit{lagt vores vigtigste punkter først på dagsordenen}）正是为了配合 Kasper 15 点接女儿。\textbf{排雷}：A 错称“取消会议”（实际只是改期到 13:30）；B 纯属无中生有。
  \item \textbf{4. B}：原文 \textit{de skal holdes, inden barnet fylder et år --- ellers bortfalder de} 与 B 中 \textit{går tabt}（作废/失效）完全同义替换。\textbf{排雷}：A 忽略 11 周 \textit{øremærket}（专款专用）不可转让；C 将“共 48 周（\textit{tilsammen}）”误作“每人单独 48 周”。
  \item \textbf{5. C}：Thomas 认为结论（\textit{konklusionen}，即报告最后部分）还可以更锐利，想明天上午 9 点一起看看（\textit{kigge på den sammen}），与 C 完全吻合。\textbf{排雷}：A 与修改需求矛盾；B 称重写引言（\textit{indledningen}），原文明确夸赞引言很强。
  \item \textbf{6. A}：原文末句指出 18 岁以下儿童及青少年逾期归还\textbf{不}收滞纳金，仅在物品遗失时（\textit{hvis materialet bliver væk}）才收赔偿账单，对应 A。\textbf{排雷}：B 忽略了前提 \textit{Hvis ingen andre har reserveret materialet}（无人预约时）；C 违背“18 岁以下不交滞纳金”。
\end{itemize}
\end{answerbox}

\begin{transbox}
\small
\textbf{\color{NCDTeal}[Opgave 4 六则短文中文对照]}\\
\textbf{1. 年假与度假金规定（borger.dk）：}作为丹麦受薪雇员，你每受雇一个月可累计 2.08 天带薪休假或度假金，折合一个完整假期年度（9 月 1 日至次年 8 月 31 日）共 25 天。你在累计满次月即可陆续使用已赚取的年假。与雇主商定休假后，你须自行记得在 borger.dk 申请领取度假金，但最早不得早于休假首日前一个月。若你在休假开始前生病，则无义务按原计划休假，可申请顺延。\\
\textbf{2. 失业保险与失业金（star.dk）：}若你失业，必须在失业第一天在 jobnet.dk 登记为求职者方可领取失业金；最早只能从在就业中心登记之日起发放。作为全职参保人，通常须加入 a-kasse 满一年且近三年内达到一定薪资收入水平才能享有全额失业金。若无正当理由自行辞职，通常会有三周等待期（期间不发放失业金）。\\
\textbf{3. 同事工作短信：}嗨卡斯珀，由于玛雅上午要见客户，周二的团队会议遗憾从 10 点改到了 13:30。我知道你平时 15 点要接女儿，所以我把我们最核心的议题排在了议程最前面。如果你到时仍须提前离席请发消息告诉我。祝好，莉内。\\
\textbf{4. 新手父母育儿假规定（borger.dk）：}在丹麦成为父母后，宝宝出生后你们总共享有 48 周带育儿津贴的育儿假。如果双方均为受薪雇员，其中各有 11 周为个人专属不可转让假期，必须在孩子满一岁前休完，否则作废。其余周数则可根据家庭与工作情况自由分配。\\
\textbf{5. 项目报告反馈短信：}嗨娜迪亚，我已经通读了你的项目报告初稿。引言和数据看起来非常扎实，但我觉得结论部分还可以再尖锐精炼一些。我们明天早上 9 点边喝咖啡边一起看一下结论部分，再发给客户好吗？祝好，托马斯。\\
\textbf{6. 公共图书馆借阅与逾期费规定（bibliotek.dk）：}在公共图书馆借阅图书、电影或桌游，借期通常为 30 天。若无人预约该资料，你最多可通过图书馆 App 或网站续借三次。到期前三天系统会自动发送短信或邮件提醒。逾期未还时，18 岁以上成人须按超期天数递增缴纳逾期费；而 18 岁以下儿童及青少年逾期不收费，但若资料遗失则须支付赔偿账单。
\end{transbox}

\paragraph*{核心高频词表（Gloser｜15 个 B1+/B2- 完整屈折词条）}

\begin{glosertable}
名词 & en arbejdsuge & \textit{(arbejdsugen, arbejdsuger, -ugerne)} & 工作周 \\
名词 & en gevinst & \textit{(gevinsten, gevinster, gevinsterne)} & 收益，好处，获益 \\
形容词 & udkørt & \textit{(udkørt, udkørt, udkørte)} & 筋疲力尽的，精疲力竭的 \\
动词 [har] & at indføre & \textit{(indfører, indførte, har indført)} & 引进，推行，实施 \\
名词 & en udfordring & \textit{(udfordringen, udfordringer, -erne)} & 挑战，难题 \\
名词 & en erhvervsuddannelse & \textit{(erhvervsuddannelsen, -uddannelser, -erne)} & 职业教育，技工培训 \\
名词 & en læreplads & \textit{(lærepladsen, lærepladser, -erne)} & 学徒岗位，实习名额 \\
名词 & et karriereskift & \textit{(karriereskiftet, karriereskift, -skiftene)} & 职业转型（单复同形） \\
名词 & en overenskomst & \textit{(overenskomsten, overenskomster, -erne)} & 劳资集体协议 \\
名词 & en barselsorlov & \textit{(barselsorloven, barselsorlover, -erne)} & 育儿假，产假/陪产假 \\
动词 [er] & at bortfalde & \textit{(bortfalder, bortfaldt, er bortfaldet)} & 失效，作废，取消（不规则！） \\
动词 [har] & at opsige & \textit{(opsiger, opsagde, har opsagt)} & 辞职；解除合同（不规则！） \\
名词 & et fællesskab & \textit{(fællesskabet, fællesskaber, -ne)} & 集体感，社群，凝聚力 \\
名词 & en ligestilling & \textit{(ligestillingen, \textnormal{通常无复数})} & 平等，平权（尤指男女平权） \\
动词 [har] & at prioritere & \textit{(prioriterer, prioriterede, har prioriteret)} & 优先考虑，排定优先顺序 \\
\end{glosertable}

% ---------------------------------------------------------------------
\subsubsection*{4.3 写作双档高分示范（Opgave 1 $\ge$ 125 ord \& Opgave 2 $\ge$ 100 ord）}

\begin{emailbox}{Skrivning Opgave 1 范文一：B1+ 稳过版（153 ord｜三问全覆 · 情态小品词自然落地）}
Kære Mikkel

\vspace{0.25em}
Tusind tak for din mail! Det er \da{jo} så dejligt at høre fra dig igen.

\vspace{0.25em}
Med vores flytning går det heldigvis rigtig godt. Vi er næsten færdige med at pakke kasserne ud, og lejligheden føles allerede hyggelig. Det nye kvarter er meget roligere end før, og der ligger både et supermarked og en god bager i nærheden, så vi er faldet godt til.

\vspace{0.25em}
Først og fremmest tillykke med forfremmelsen til teamleder! Det har du \da{da} virkelig fortjent. \dabold{Hvis jeg var dig, ville jeg vælge} en ekstra uges ferie frem for mere i løn. Man betaler \da{jo} alligevel meget i skat af en lønforhøjelse, og mere tid med familien giver langt mere energi i længden.

\vspace{0.25em}
Hvad angår arbejdsmails i weekenden, synes jeg ikke, at det er en god idé. Alle har brug for rigtig fri for at undgå stress, \dabold{medmindre} der sker noget akut.

\vspace{0.25em}
Skriv snart, hvad du beslutter!

\vspace{0.25em}
Kærlig hilsen\\
Linnea

\vspace{0.3em}
\rule{\linewidth}{0.4pt}\\[0.2em]
\small\textbf{\color{NCDTeal}[考官点评 · 153 ord]}：稳居 125 词以上安全区。依次回应 Mikkel 全部三项提问（搬家安置与新社区适应、升职选休假还是加薪、周末办公邮件文化），恰当嵌入 \textit{jo / da} 表达熟人共情，并用第二条件句（\textit{Hvis jeg var dig, ville jeg...}）与限制条件词 \textit{medmindre}，逻辑清晰稳过。
\end{emailbox}

\begin{emailbox}{Skrivning Opgave 1 范文二：B2- 高分冲刺版（192 ord｜假拟条件 + 被动语态 + 让步反驳 + 情态小品词）}
Kære Mikkel

\vspace{0.25em}
Hvor er det dejligt at høre fra dig! Undskyld, at jeg ikke har skrevet før, men flytningen har \da{jo} taget al vores tid de seneste uger.

\vspace{0.25em}
Vi er nu kommet på plads i den nye lejlighed og er allerede faldet utrolig godt til. Kvarteret er grønnere, end jeg havde forventet, og naboen har ovenikøbet inviteret os på kaffe. Det var \da{egentlig} lidt hårdt at flytte væk fra vores gamle venner, men vi fortryder ikke skiftet.

\vspace{0.25em}
Et kæmpe tillykke med, at du \dabold{er blevet forfremmet} til teamleder! \dabold{Hvis jeg stod i dit sted, ville jeg} uden tvivl vælge den ekstra uges ferie. \dabold{Ganske vist} lyder en højere månedsløn fristende, \dabold{men} efter skat er forskellen \da{jo} ikke så stor. Mere ferie giver dig derimod tid til at koble helt af.

\vspace{0.25em}
Med hensyn til arbejdsmails i weekenden mener jeg, at en sådan kultur \dabold{bør bremses} i tide. Når ledere skriver lørdag aften, bliver medarbejderne usikre på, om de \dabold{forventes} at svare straks. Arbejde og fritid \dabold{må adskilles}, \dabold{medmindre} en akut sag kræver handling.

\vspace{0.25em}
Held og lykke! Det skal \da{nok} gå rigtig godt i den nye rolle! Lad os snart mødes!

\vspace{0.25em}
Bedste hilsner\\
Linnea

\vspace{0.3em}
\rule{\linewidth}{0.4pt}\\[0.2em]
\small\textbf{\color{NCDNavy}[高分亮点 · 192 ord]}：\ding{172} \textit{blive} 被动（\textit{er blevet forfremmet}）与 \textit{-s} 规范被动（\textit{bør bremses, forventes, må adskilles}）自然切换；\ding{173} 高阶假拟条件（\textit{Hvis jeg stod i dit sted, ville jeg...}）搭配让步反驳（\textit{Ganske vist..., men...}）；\ding{174} 情态小品词 \textit{jo / egentlig / nok} 精准拉近私人信件亲密度。
\end{emailbox}

\begin{emailbox}{Skrivning Opgave 2 范文一：B1+ 稳过版（正文 131 ord / 含标题 136 ord｜正反双面权衡）}
\textbf{\color{NCDNavy}Bevar både friheden og fællesskabet}

\vspace{0.25em}
Jeg forstår godt, hvorfor Henrik Lind er bekymret for fællesskabet på kontoret, men jeg er ikke enig i hans løsning.

\vspace{0.25em}
På den ene side har han ret i, at det er vigtigt at mødes fysisk. Når man sidder sammen, er det nemmere at hjælpe nye kollegaer, løse problemer hurtigt og føle sig som en del af et hold. Hvis kontoret altid er tomt, bliver arbejdsdagen ensom.

\vspace{0.25em}
På den anden side giver hjemmearbejde også mange fordele i hverdagen. Mange sparer lang transporttid, koncentrerer sig bedre om svære opgaver og får en sundere balance mellem arbejde og familieliv.

\vspace{0.25em}
Derfor bør virksomheder ikke tvinge folk ind på kontoret hver dag. I stedet kan man aftale faste fællesdage, hvor hele teamet mødes to eller tre gange om ugen. Så bevarer man både friheden og sammenholdet.

\vspace{0.3em}
\rule{\linewidth}{0.4pt}\\[0.2em]
\small\textbf{\color{NCDTeal}[考官点评 · 131/136 ord]}：远超 100 词要求。开篇精准表态（理解担忧但不同意强制全勤），中段采用 \textit{På den ene side / På den anden side} 严密兼顾团队互助与通勤减压，末段提出“每周固定两三天共同到岗”的可行折中方案，紧扣 Modul 4 辩论规范。
\end{emailbox}

\begin{emailbox}{Skrivning Opgave 2 范文二：B2- 高分冲刺版（正文 159 ord / 含标题 165 ord｜让步反驳 + 假拟条件 + \textit{-s} 被动）}
\textbf{\color{NCDNavy}Hybridarbejde styrker både trivsel og sammenhold}

\vspace{0.25em}
Henrik Lind rejser en vigtig debat om arbejdsmiljø, men jeg er grundlæggende uenig i hans konklusion om helt at afskaffe hjemmearbejdet.

\vspace{0.25em}
\dabold{Ganske vist} kan for meget arbejde bag hjemmeskærmen svække den uformelle kontakt ved kaffemaskinen og gøre det sværere for nyansatte at blive en naturlig del af arbejdspladsens kultur. \dabold{Ikke desto mindre} overser Lind de store gevinster, som fleksible rammer skaber. Når medarbejdere slipper for spildtid i myldretiden og kan fordybe sig i ro derhjemme, falder både stressniveauet og sygefraværet ofte mærkbart.

\vspace{0.25em}
\dabold{Hvis medarbejdere blev tvunget} ind på kontoret fem dage om ugen, \dabold{ville} mange erfarne folk sandsynligvis søge videre til mere moderne virksomheder. Et stærkt fællesskab \dabold{skabes} nemlig ikke gennem tvungent fremmøde, men gennem god ledelse, fælles mål og tydelige rammer.

\vspace{0.25em}
Efter min mening ligger løsningen i en hybridmodel med to faste kontordage om ugen, hvor møder og tæt samarbejde \dabold{prioriteres}, mens opgaver, der kræver dyb ro, \dabold{løses} hjemme. På den måde \dabold{sikres} både trivsel og fagligt sammenhold.

\vspace{0.3em}
\rule{\linewidth}{0.4pt}\\[0.2em]
\small\textbf{\color{NCDNavy}[高分亮点 · 159/165 ord]}：\ding{172} 标准报刊评论让步架构 \textit{Ganske vist..., Ikke desto mindre...}；\ding{173} 假拟条件嵌套 \textit{blive} 被动（\textit{Hvis medarbejdere blev tvunget..., ville... søge videre}）论证强制到岗的反向人才流失后果；\ding{174} 连续精准运用机制性 \textit{-s} 被动（\textit{skabes, prioriteres, løses, sikres}），句式紧凑老练。
\end{emailbox}

% ---------------------------------------------------------------------
\subsubsection*{4.4 口语示范逐字稿、核心 IPA 与双人社会议题辩论}

\begin{transbox}
\small
\textbf{\color{NCDNavy}[Emneopgave 1A 口语陈述示范：Arbejdsliv: balance mellem job og fritid｜165 ord · 约 1.5 分钟]}\\
\textit{Jeg vil gerne fortælle om arbejdsliv i Danmark med fokus på balancen mellem job og fritid \ipa{ˈfʁiːˌtˢiðˀ}. På det danske arbejdsmarked arbejder en fuldtidsansat normalt 37 timer om ugen, og der lægges stor vægt på et sundt arbejdsmiljø og et fladt hierarki \ipa{hiaˈkʰiːˀ}. Det er for eksempel helt almindeligt, at både mænd og kvinder går klokken fire for at hente børn eller dyrke motion, fordi ledelsen bygger på tillid \ipa{ˈtˢilˌliːˀð} frem for kontrol. Fordelen ved denne fleksibilitet er, at medarbejderne ofte får mere overskud og mindre stress i hverdagen. Udfordringen \ipa{ˈuðˌfɒːˀdʁeŋˀən} er dog, at grænsen mellem arbejde og privatliv nogle gange bliver uklar, når man kan tjekke mails på telefonen om aftenen. I mit eget arbejdsliv prøver jeg derfor at holde klare rammer: jeg arbejder effektivt i dagtimerne, men jeg slukker arbejdsmailen klokken fem og prioriterer løbeture og madlavning med familien. Hvis jeg kunne ændre én ting i fremtiden, ville jeg gerne prøve en fire dages arbejdsuge, fordi tre fridage ville give endnu bedre tid til hvile og fritid.}\\[4pt]
\textbf{\color{NCDTeal}[参考译文]}：我想聊聊丹麦的职场生活，重点是工作与休闲之间的平衡。在丹麦劳动力市场上，全职员工通常每周工作 37 小时，而且非常注重健康的工作环境与扁平的层级文化。例如，无论是男性还是女性，四点下班去接孩子或运动都非常普遍，因为管理靠的是信任而非监控。这种灵活性的优点在于员工在日常生活中往往更有精力、压力更小。不过挑战在于，当晚上还能在手机上查收邮件时，工作与私生活之间的界限有时会变得模糊。因此在我自己的职场生活中，我尽量保持清晰的界线：白天高效工作，但五点就关掉工作邮箱，把时间优先留给跑步和与家人一起做饭。如果将来能改变一件事，我很想尝试四天工作制，因为三天休息日能提供更加充裕的休息和休闲时间。
\end{transbox}

\paragraph*{Emneopgave 1A Del 2 考官追问与高分应答示范（4 题）：}
\dialogturn{1}{Eksaminator}{Vil du fortælle lidt mere om, hvordan du selv oplever det flade hierarki og tilliden på en dansk arbejdsplads?}{hierarki \ipa{hiaˈkʰiːˀ} · tilliden \ipa{ˈtˢilˌliːˀðən}}{你能再多谈谈自己亲身感受到的丹麦职场扁平层级和信任文化吗？}
\dialogturn{2}{Kursist}{Ja, jeg oplever det ved, at vores leder ikke sidder og tjekker, hvornår vi møder præcist, så længe opgaverne bliver løst til tiden. Man kan også altid sige sin ærlige mening til møder uden at være bange for chefen, hvilket skaber et trygt arbejdsmiljø.}{opgaverne \ipa{ˈʌbˌɡæːvɐnə} · trygt \ipa{ˈtˢʁyɡd}}{好的，我的感受是只要任务按时完成，领导不会盯着看我们精确几分钟打卡。而且在会议上大家随时可以坦诚发表意见而不用惧怕老板，这营造出非常有安全感的工作氛围。}
\dialogturn{3}{Eksaminator}{Kan du give nogle eksempler på, hvordan du sætter grænser mellem arbejde og fritid i din egen hverdag?}{grænser \ipa{ˈɡʁɛnsɐ} · hverdag \ipa{ˈvɛːɐˌdæːˀ}}{你能举几个例子说明你在日常生活中如何划清工作与休闲的界限吗？}
\dialogturn{4}{Kursist}{For det første har jeg slået notifikationer fra arbejdsmailen fra på min private telefon. For det andet aftaler jeg faste tider til sport eller middag med venner efter fyraften, så jeg ikke bliver siddende ved computeren hele aftenen.}{notifikationer \ipa{notˢifikaˈɕoːˀnɐ} · fyraften \ipa{ˈfyːɐˌɑfdən}}{首先，我在私人手机上关闭了工作邮箱的消息推送。其次，下班后我会固定约好运动或和朋友吃晚饭的时间，这样就不会整个晚上都黏在电脑前。}
\dialogturn{5}{Eksaminator}{Hvad er efter din mening fordelene og ulemperne ved at arbejde hjemmefra et par dage om ugen?}{fordelene \ipa{ˈfɒːˌdeːˀlənə} · ulemperne \ipa{ˈuːˌlɛmbɐnə}}{在你看来，每周居家办公几天有哪些优点和缺点？}
\dialogturn{6}{Kursist}{Fordelen er helt klart, at man sparer transporttid og kan arbejde meget koncentreret uden støj. Ulempen er derimod, at man kan komme til at savne den uformelle snak med kollegaerne, og det kan være sværere for nye medarbejdere at lære arbejdspladsen at kende.}{koncentreret \ipa{kʰʌnsənˈtˢʁeːˀɐð} · uformelle \ipa{ˈuːfʌˌmɛlˀə}}{优点很明显，就是省去通勤时间且能在无噪音环境中高度专注。但反过来说，缺点是容易想念和同事的随性闲聊，而且新员工也更难熟悉公司文化。}
\dialogturn{7}{Eksaminator}{Hvad synes du om forskellen mellem arbejdskulturen i Danmark og i dit hjemland? Hvorfor?}{arbejdskulturen \ipa{ˈɑːbɑi̯dsˌkʰulˌtˢuːˀɐn} · hjemland \ipa{ˈjɛmˌlanˀ}}{你对丹麦和你母国工作文化的差异怎么看？为什么？}
\dialogturn{8}{Kursist}{I mit hjemland er arbejdstiden ofte længere, og det ses tit som et tegn på flid, hvis man bliver sent på kontoret. I Danmark handler det derimod om effektivitet og familiebalance, og det synes jeg faktisk giver både gladere mennesker og et mere bæredygtigt arbejdsliv.}{effektivitet \ipa{efɛɡtˢiviˈtˢeːˀd} · bæredygtigt \ipa{ˈbɛːɐˌdyɡdiɡd}}{在我母国工作时间通常更长，待在办公室加班常被视为勤奋的表现。而在丹麦更看重效率与家庭平衡，我其实觉得这样既让员工更快乐，也让职场生涯更可持续。}

\begin{transbox}
\small
\textbf{\color{NCDNavy}[Emneopgave 1B 口语陈述示范：Ligestilling: barsel og karriere｜163 ord · 约 1.5 分钟]}\\
\textit{Jeg har valgt at fortælle om ligestilling \ipa{ˈliːəˌsdeleŋ} i Danmark med fokus på barsel og karriere \ipa{kʰɑˈjɛːɐ}. I Danmark har nye forældre tilsammen ret til 48 ugers barselsorlov med barselsdagpenge, og siden 2022 har 11 af ugerne været øremærket \ipa{ˈøːɐˌmɛɐ̯ɡəð} til hver forælder, hvis begge er lønmodtagere. Formålet med disse regler er at skabe mere ligestilling både i hjemmet og på arbejdsmarkedet. Når fædre også tager en længere periode på barsel, får de et tættere forhold til barnet, og samtidig bliver det ikke kun kvinderne, der sætter karrieren på pause. Alligevel findes der stadig udfordringer, fordi nogle fædre frygter, at langt fravær kan påvirke deres løn eller muligheder for forfremmelse \ipa{fʌˈfʁɑmˀəlsə} i private firmaer. I min egen omgangskreds har jeg dog set flere danske fædre nyde deres barsel i fulde drag, hvilket er meget anderledes end i mit hjemland, hvor mænd sjældent tager mere end et par ugers fri. Samlet set synes jeg, at øremærket barsel er et rigtig vigtigt skridt mod reel ligestilling i vores samfund.}\\[4pt]
\textbf{\color{NCDTeal}[参考译文]}：我选择谈论丹麦的平权问题，重点是育儿假与职业发展。在丹麦，新手父母合计享有 48 周带育儿津贴的育儿假，自 2022 年起若双方均为受薪雇员，每人各有 11 周专属不可转让假期。这些规则的目的是在家庭内部和劳动力市场上创造更多平等。当父亲也休较长育儿假时，不仅能与孩子建立更亲密的关系，同时也避免只有女性不得不暂停职业发展。尽管如此，挑战依然存在，因为部分父亲担心长期缺勤会影响自己在私营企业中的薪酬或晋升机会。不过在我自己的朋友圈里，我看到好几位丹麦爸爸尽情享受他们的育儿假时光，这与我的母国截然不同——在那里男性很少休假超过一两周。总体而言，我认为专属育儿假是迈向社会真正平权的重要一步。
\end{transbox}

\paragraph*{Emneopgave 1B Del 2 考官追问与高分应答示范（4 题）：}
\dialogturn{1}{Eksaminator}{Vil du fortælle lidt mere om, hvorfor du mener, at øremærket barsel til fædre er vigtig?}{øremærket \ipa{ˈøːɐˌmɛɐ̯ɡəð} · vigtig \ipa{ˈveɡdi}}{你能再多讲讲为什么你认为专门留给父亲的育儿假很重要吗？}
\dialogturn{2}{Kursist}{Fordi ændringer i loven sender et stærkt signal til både familier og arbejdsgivere om, at omsorg for børn er et fælles ansvar. Når det bliver normalt for mænd at gå på barsel i tre måneder, bliver det lettere for den enkelte far at sige ja uden dårlig samvittighed over for kollegaerne.}{omsorg \ipa{ˈʌmˌsɒːˀ} · samvittighed \ipa{sɑmˈvidihəðˀ}}{因为法律层面的调整向家庭和雇主都发出了强有力的信号，表明育儿是共同责任。当男性休三个月育儿假成为常态，个别父亲也更容易坦然接受而不会对同事心怀愧疚。}
\dialogturn{3}{Eksaminator}{Kan du give nogle eksempler på, hvordan arbejdsdelingen i familien kan påvirke kvinders og mænds karriere?}{arbejdsdelingen \ipa{ˈɑːbɑi̯dsˌdeːleŋən} · påvirke \ipa{ˈpʰʌˌviɐ̯ɡə}}{你能举例说明家庭内部分工如何影响女性和男性的职业生涯吗？}
\dialogturn{4}{Kursist}{Ja, hvis det altid er moderen, der tager barnets første sygedag og henter tidligt fra vuggestue, går hun ofte ned i tid og sakker bagud med løn og pension. Hvis forældrene derimod skiftes, kan begge to fortsætte med at udvikle sig fagligt.}{vuggestue \ipa{ˈvoɡəˌsduːə} · sakker bagud \ipa{ˌsɑɡɐ ˈbæːˀˌuðˀ}}{是的，如果总是母亲在孩子生病第一天请假、早早去托儿所接娃，她往往就会转成兼职并在薪资和养老金上落后。相反，如果父母轮流承担，双方都能继续在专业上成长。}
\dialogturn{5}{Eksaminator}{Hvad er fordelene og ulemperne ved, at staten bestemmer, at 11 uger ikke kan overføres til partneren?}{bestemmer \ipa{beˈsdɛmˀɐ} · overføres \ipa{ˈʌu̯ɐˌføːˀɐs}}{政府规定其中 11 周不能转给伴侣，这有哪些利弊？}
\dialogturn{6}{Kursist}{Fordelen er, at faren bliver motiveret til at tage sin del af orloven, hvilket styrker ligestillingen på sigt. Ulempen er mindre fleksibilitet for de familier, hvor faren for eksempel er selvstændig med en lille butik, eller hvor begge hellere selv ville planlægge ugerne frit.}{motiveret \ipa{motˢiˈveːˀɐð} · selvstændig \ipa{ˈsɛlˌsdɛnˀdi}}{优点是能激励父亲承担自己那份假期，从长远看促进平权。缺点是对一些家庭来说缺乏灵活度，比如父亲若是开小店的个体户，或夫妻俩本想更自由地自行分配周数。}
\dialogturn{7}{Eksaminator}{Hvad synes du om barselsreglerne i Danmark sammenlignet med reglerne i dit hjemland? Hvorfor?}{sammenlignet \ipa{ˈsɑmənˌliːˀnəð} · reglerne \ipa{ˈʁɛːˀlɐnə}}{与你母国的规定相比，你如何评价丹麦的育儿假制度？为什么？}
\dialogturn{8}{Kursist}{Jeg synes, den danske model er meget generøs og familievenlig, fordi både mor og far er beskyttet økonomisk gennem barselsdagpenge. I mit hjemland er orloven kortere, især for mænd, så det danske system gør det langt nemmere at kombinere børn og karriere.}{generøs \ipa{ɕenəˈʁøːˀs} · beskyttet \ipa{beˈsɡydəð}}{我觉得丹麦的模式非常慷慨且对家庭友好，因为父母双方在经济上都受到育儿津贴的保障。在我母国假期较短，尤其是男性几乎没多少假，因此丹麦体制让兼顾育儿与事业容易得多。}

\vspace{0.35em}
\hrule
\vspace{0.35em}
\paragraph*{Mundtlig Opgave 2 Del 1（个人看图必答三题示范 · Kursist A \& Kursist B）：}

\begin{grammatikbox}{Kursist A 个人应答示范（图片 A：居家远程办公｜含 IPA 与译文）}
\small
\begin{enumerate}
  \item \textbf{Hvorfor tror du, de arbejder hjemmefra?} \textit{Jeg tror, mange vælger at arbejde hjemmefra \ipa{ˈjɛməˌfʁɑːˀ}, fordi det sparer en masse transporttid i myldretiden. Desuden giver det en rolig atmosfære, hvor man kan koncentrere sig uden at blive afbrudt af kollegaer i et stort kontorlandskab.}（我认为许多人选择居家办公是因为能省下高峰期的大量通勤时间，而且家里安静，不会被开放式办公室的同事打断，更能专注工作。）
  \item \textbf{Hvilke fordele kan der være ved det?} \textit{Den største fordel er fleksibiliteten \ipa{flɛɡsibiliˈtˢeːˀdən} i hverdagen. Man kan sove lidt længere, sætte en tøjvask over i frokostpausen og få familielivet til at hænge bedre sammen. Samtidig oplever mange mindre stress og bedre fordybelse \ipa{fʌˈdyːˀbəlsə}.}（最大优点在于日常的灵活性：能多睡一会儿、午休时顺手洗锅衣服，家庭生活更容易统筹；同时很多人感到压力更小、工作更沉浸。）
  \item \textbf{Mener du, der kan være nogen ulemper ved at arbejde hjemmefra? Hvorfor? / Hvorfor ikke?} \textit{Ja, jeg mener helt klart, at der også er ulemper. For det første kan man hurtigt føle sig isoleret \ipa{isoˈleːˀɐð} og savne den sociale kontakt med sine kollegaer. For det andet kan det være svært at slukke arbejdet mentalt, når stuen også bliver til et kontor, og nogle får endda ondt i ryggen, hvis de ikke har et ordentligt skrivebord.}（是的，我认为显然也有缺点。首先人很容易感到孤立并想念与同事的社交互动；其次当客厅变成办公室时，精神上很难“下班”，没有像样书桌还可能腰背酸痛。）
\end{enumerate}
\end{grammatikbox}

\begin{grammatikbox}{Kursist B 个人应答示范（图片 B：兼职父母提前接娃｜含 IPA 与译文）}
\small
\begin{enumerate}
  \item \textbf{Hvorfor tror du, de vælger deltid?} \textit{Jeg tror, nogle forældre vælger deltid \ipa{ˈdeːlˌtˢiðˀ}, mens børnene er små, fordi de ønsker mere nærvær og ro i familien. Når begge forældre har fuldtidsarbejde, kan eftermiddagene føles meget stressede med indkøb, madlavning og trætte børn.}（我认为一些父母在孩子年幼时选择兼职，是因为希望多些陪伴与家庭的从容。若父母双双全职，下午买菜、做饭加照顾疲惫的孩子会让人觉得倍感压力。）
  \item \textbf{Hvilke fordele kan der være ved det?} \textit{Fordelen er, at barnet får kortere dage i børnehaven og tid til at lege i ro og mag derhjemme. For forældrene giver det mindre dårlig samvittighed \ipa{sɑmˈvidihəðˀ} og bedre overskud til både praktiske pligter og hyggelige stunder med familien.}（优点在于孩子在幼儿园的时间更短，回家有充裕时间安安静静玩耍。对父母而言能减少愧疚感，也更有精力应付日常家务与亲子温馨时光。）
  \item \textbf{Mener du, der kan være nogen ulemper ved at arbejde på deltid? Hvorfor? / Hvorfor ikke?} \textit{Ja, der er tydelige ulemper på længere sigt. Først og fremmest tjener man færre penge her og nu og sparer også mindre op til sin pension \ipa{pʰɑŋˈɕoːˀn}. Derudover er det desværre stadig oftest kvinder, der går ned i tid, hvilket kan bremse deres karriereudvikling og skabe ulighed i parforholdet.}（是的，从长远来看有明显缺点。首先眼下收入更少，养老金积累也随之减少；此外遗憾的是目前仍多为女性选择缩减工时，这可能阻碍她们的职业晋升并在伴侣关系中造成不平等。）
\end{enumerate}
\end{grammatikbox}

\vspace{0.35em}
\hrule
\vspace{0.35em}
\paragraph*{Mundtlig Opgave 2 Del 2（Paropgave：四天带薪工作制 · 8 轮双人辩论示范，兼顾个人与社会维度）：}

\dialogturn{1}{Kursist A}{Det er et virkelig spændende emne! Personligt synes jeg, at en fire dages arbejdsuge med fuld løn lyder fantastisk for den enkelte medarbejder. Man får jo mere tid til familien og hvile. Hvad tænker du om det?}{arbejdsuge \ipa{ˈɑːbɑi̯dsˌuːə} · personligt \ipa{pʰɛɐ̯ˈsoːˀnlið}}{这是个非常有趣的话题！我个人觉得，带全薪的四天工作制对员工个人来说听起来太棒了，毕竟能有更多时间陪伴家人和休息。你怎么看？}
\dialogturn{2}{Kursist B}{Jeg er enig i, at det lyder tillokkende for den enkelte, fordi færre arbejdsdage kan mindske stress og give højere livskvalitet. Men jeg spekulerer på, om arbejdsdagene ikke bliver alt for lange og intense, hvis man skal nå præcis det samme på fire dage?}{tillokkende \ipa{ˈtˢeˌlʌɡənə} · livskvalitet \ipa{ˈliːu̯sˌkʰvaliˌtˢeːˀd}}{我同意这对个人很有吸引力，因为工作日变少能减轻压力、提高生活质量。但我在想，如果四天内要完成完全相同的工作量，每天会不会变得太长太紧张？}
\dialogturn{3}{Kursist A}{Det er en god pointe. Hvis opgaverne ikke skæres ned, risikerer man bare mere overarbejde. Men hvad mener du egentlig om fordelene for samfundet som helhed? Kan det ikke sænke sygefraværet og udgifterne i sundhedsvæsenet?}{sygefraværet \ipa{ˈsyːəˌfʁɑvɛːˀɐð} · sundhedsvæsenet \ipa{ˈsɔnhəðsˌvɛːˀsnəð}}{说得好。如果任务量不减，确实只会导致更多加班。不过你对全社会层面的好处怎么看？这难道不能降低病假率并减少医疗体系支出吗？}
\dialogturn{4}{Kursist B}{Jo, på lang sigt kan et lavere stressniveau helt klart spare samfundet for penge til behandling af udbrændthed. Samtidig har borgerne også mere tid til frivilligt arbejde i foreninger. Men hvordan ser du ulemperne for samfundet, for eksempel på hospitaler og i skoler?}{udbrændthed \ipa{ˈuðˌbʁɛndˀheːð} · frivilligt \ipa{ˈfʁiˌvilið}}{没错，从长远看更低的压力水平显然能帮社会省下治疗职业倦怠的开支，而且市民也有更多时间参与社团志愿活动。不过你怎么看它对社会的弊端，比如医院和学校？}
\dialogturn{5}{Kursist A}{Ja, der ligger nok den største ulempe for samfundet. I velfærdsfag som plejehjem, børnehaver og sygehuse kan man ikke bare arbejde hurtigere --- der skal jo være personale hele tiden. Hvis alle kun arbejdede fire dage med fuld løn, ville det blive meget dyrt for staten, medmindre man ansatte flere folk.}{velfærdsfag \ipa{ˈvɛlˌfɛːɐ̯sˌfæːˀ} · personale \ipa{pʰɛɐ̯soˈnæːlə}}{是的，这大概就是对社会最大的弊端所在。在养老院、幼儿园和医院这类福利行业，工作没法靠“干快点”解决——必须时刻有人在岗。如果所有人都上四天领全薪，除非增聘人手，否则国家财政负担会极其沉重。}
\dialogturn{6}{Kursist B}{Præcis! Og i Danmark har vi allerede mangel på arbejdskraft i netop de brancher. Så hvis den offentlige sektor ikke kan tilbyde det samme som private IT-firmaer, risikerer vi jo endnu større ulighed på arbejdsmarkedet. Tror du, der findes en realistisk løsning?}{arbejdskraft \ipa{ˈɑːbɑi̯dsˌkʰʁɑfd} · ulighed \ipa{ˈuːˌliːheːðˀ}}{正是如此！而且丹麦目前恰好就在这些行业闹用工荒。所以如果公共部门无法提供与私营 IT 公司同等的待遇，劳动力市场就会面临更大的不平等。你觉得有切实可行的解决办法吗？}
\dialogturn{7}{Kursist A}{Måske bør man ikke tvinge én fast model ned over hele samfundet, men lade hver branche vælge selv. På et kontor kan fire dages uge fungere godt, mens man i sundhedsvæsenet måske kan give højere løn eller mere fleksible vagter i stedet. Hvad synes du om den idé?}{branche \ipa{ˈbʁɑŋɕə} · fleksible \ipa{flɛɡˈsiːblə}}{也许不应强迫整个社会套用单一模式，而是由各行业自行选择。办公室岗位适合推四天工作制，而医疗领域或许可以通过提高薪酬或提供更灵活排班作为替代。你觉得这个主意怎么样？}
\dialogturn{8}{Kursist B}{Det synes jeg lyder meget fornuftigt! På den måde får den enkelte mere frihed, hvor det er muligt, uden at velfærdssamfundet mister sin sammenhængskraft.}{fornuftigt \ipa{fʌˈnʌfdiɡd} · sammenhængskraft \ipa{ˈsɑmənˌhɛŋsˌkʰʁɑfd}}{我觉得这听起来非常合理！这样既能在可行岗位上赋予个人更多自由，又不会让福利社会丧失凝聚力。}
